\begin{table}[t]
\centering
\caption{Standardized Effect Sizes}
\label{tab:sde}
\begin{tabular}{lcccccc}
\hline\hline
Outcome & $\hat{\beta}$ & SE & SD($Y$) & SDE & SE(SDE) & Classification \\
\hline
Mean price & -0.597 & 0.288 & 1.981 & -0.301 & 0.145 & Large negative \\
Log price & -0.042 & 0.025 & 0.137 & -0.309 & 0.183 & Large negative \\
Peak price & -1.022 & 0.442 & 2.172 & -0.470 & 0.204 & Large negative \\
Off-peak price & -0.913 & 0.414 & 2.082 & -0.438 & 0.199 & Large negative \\
\hline\hline
\end{tabular}
\begin{tablenotes}
\item \textit{Notes:} \textbf{Country:} Japan. \textbf{Research question:} Do post-Fukushima nuclear reactor restarts, approved through Japan's Nuclear Regulation Authority safety review process, reduce wholesale electricity prices in the affected regional grid? \textbf{Policy mechanism:} NRA-approved reactor restarts add low-marginal-cost nuclear baseload generation to regional electricity grids, potentially displacing high-cost LNG-fired generation from the merit order and reducing the market-clearing price in the Japan Electric Power Exchange day-ahead auction. \textbf{Outcome definition:} Monthly mean (or log, peak, off-peak) of half-hourly JEPX day-ahead spot auction settlement prices in Yen per kilowatt-hour at the regional area level. \textbf{Treatment:} Binary indicator for whether any reactor in a region has restarted (first restart ranges from August 2015 to December 2024 across five treated regions). \textbf{Data:} JEPX day-ahead spot prices from japanesepower.org, 9 regions, January 2012 to December 2025, aggregated to 1,512 region-months from 2.5 million half-hourly observations. \textbf{Method:} Callaway and Sant'Anna (2021) staggered DiD with not-yet-treated control group; standard errors via multiplier bootstrap. \textbf{Sample:} Nine JEPX electricity regions (excluding Okinawa); five treated regions with NRA-approved restarts and four never-treated control regions; Japan's 50Hz/60Hz frequency split limits cross-regional price spillovers. SDE $= \hat{\beta} / \text{SD}(Y)$ where SD($Y$) is the pre-treatment standard deviation. Classification refers to magnitude, not statistical significance: Large ($|$SDE$| > 0.15$), Moderate ($0.05$--$0.15$), Small ($0.005$--$0.05$), Null ($< 0.005$).
\end{tablenotes}
\end{table}

